\section{Introduction}

Can you identify a language model's architecture family from its adversarial failures alone ---
without knowing its weights, training data, or any other metadata? We show the answer is yes,
in principle. Architecture family membership (encoder-only, decoder-only, encoder-decoder) accounts
for approximately 29\% of variance in normalized adversarial vulnerability across attack categories
at base model scale (permutation MANOVA $\eta^2=0.293$, Cohen's $f^2 \approx 0.41$), and this effect
is reproducible across 83\% of the adversarial attack types we evaluate. The fingerprint is real
and large. But statistically confirming it requires a larger model pool than the nine models we
study here --- a limitation we quantify precisely.

The practical stakes are concrete. A security auditor selecting between a BERT-style encoder-only
model and a GPT-style decoder-only model for an adversarial-critical deployment cannot currently
quantify their different adversarial risk profiles from first principles. Robustness assessments
today are model-specific and non-transferable: every new model requires a fresh evaluation from
scratch. Architecture-family-level fingerprinting, if confirmed, would enable auditors to make
principled inferences about an unseen model's likely vulnerability profile from its architecture
family membership alone --- a substantial reduction in evaluation cost.

The core problem runs deeper than evaluation efficiency. Existing adversarial robustness studies
evaluate models as individuals, reporting scalar accuracy-drop metrics aggregated across attack types.
This scalar aggregation discards the multivariate structure of adversarial vulnerability: different
attack types stress different model components, and whether a model is more vulnerable to lexical
substitutions than to paraphrase attacks is systematically related to how its attention mechanism
processes input perturbations. Encoder-only models with bidirectional attention globally redistribute
attention mass toward perturbed tokens; decoder-only models with causal attention process perturbations
locally in sequence order. These structural differences, we hypothesize, should produce characteristic
per-attack-category vulnerability profiles --- a $\Delta^*$-vector --- that are consistent within
architecture families and discriminable between them.

Prior work takes steps in this direction. \citet{Wang2021Adversarial} establish AdvGLUE as a
multi-category adversarial benchmark covering five GLUE tasks and multiple attack methods.
\citet{Nie2020Adversarial} provide ANLI as a human-crafted adversarial NLI benchmark.
\citet{Neerudu2023Robustness} observe that GPT-2 is more robust than BERT and T5 on GLUE
perturbations, providing a directional signal consistent with architecture-family effects. Yet none
of these works quantifies the effect size of between-family differences, formalizes the
$\Delta^*$-vector representation, or applies multivariate testing suited to small model pools.
\citet{Sun2024TrustLLM} evaluate 16 LLMs across six trustworthiness dimensions but do not
stratify by architecture family. \citet{Otmakhova2025FLUKE} demonstrate that a model's ability
to use a linguistic feature does not imply robustness to perturbations of that feature ---
consistent with our motivation but orthogonal in method. The gap is precise: no prior work
measures between-family $\Delta^*$-vector effect size under scale-matched, multi-attack-category,
statistically controlled conditions.

We address this gap with a four-part contribution. First, we introduce the \textbf{$\Delta^*$-vector
framework} for multivariate architecture-family adversarial profiling. Each model is represented
as a vector of normalized per-category vulnerability scores
($\Delta^* = (\text{Acc}_\text{clean} - \text{Acc}_\text{adv})/\text{Acc}_\text{clean}$), and
architecture-family effects are tested with permutation MANOVA --- a non-parametric method
appropriate for small $N$. Second, we provide the \textbf{first formal effect-size measurement}
of between-family $\Delta^*$ clustering: $\eta^2=0.293$ across nine models and six attack
categories, exceeding the pre-specified $\eta^2=0.15$ threshold in five of six categories (83\%).
Third, we contribute a \textbf{validated seven-module pipeline} (1,788 lines, Python) covering
fine-tuning, adversarial evaluation, $\Delta^*$ computation, statistical analysis, and
visualization. Fourth, we provide a \textbf{power analysis} establishing $N \geq 15$ ($\geq 5$
models per family) as the minimum for 80\% statistical power at the observed effect size, giving
concrete design guidance for the confirmatory follow-up study.

The paper proceeds as follows. Section~\ref{sec:related} reviews adversarial robustness benchmarks,
architecture comparison studies, and trustworthiness evaluation frameworks. Section~\ref{sec:method}
describes the $\Delta^*$-vector framework and statistical methods. Section~\ref{sec:setup} presents
experimental setup. Section~\ref{sec:results} reports results. Section~\ref{sec:discussion} discusses
implications and limitations. Section~\ref{sec:conclusion} concludes.
